% Preámbulo común de los documentos de Nexus: Normas APA 7.a edición (trabajo de estudiante).
% Se carga con % Preámbulo común de los documentos de Nexus: Normas APA 7.a edición (trabajo de estudiante).
% Se carga con % Preámbulo común de los documentos de Nexus: Normas APA 7.a edición (trabajo de estudiante).
% Se carga con % Preámbulo común de los documentos de Nexus: Normas APA 7.a edición (trabajo de estudiante).
% Se carga con \input{../comun/preambulo} justo después de \documentclass.
\usepackage[spanish,es-noshorthands,es-tabla]{babel}
\usepackage[T1]{fontenc}
\usepackage{newtxtext,newtxmath}
\usepackage{csquotes}
\usepackage[style=apa,backend=biber]{biblatex}
\addbibresource{../comun/referencias.bib}
\usepackage{graphicx}
\usepackage{booktabs}
\usepackage{tabularx}
\usepackage{array}
\usepackage{enumitem}
\usepackage{xurl}
\usepackage{tikz}
\usetikzlibrary{positioning,arrows.meta,shapes.misc}
\hypersetup{hidelinks}
\setlist{itemsep=2pt,topsep=4pt,parsep=0pt}
\setlength{\emergencystretch}{3em}

% Datos de la portada
\authorsnames{Neyder Sebastián Orozco Villamil, Deisy Lorena Jiménez Martínez}
\authorsaffiliations{Universidad Central}
\course{Modelado 3D y Video Juegos (43391312)}
\professor{Mg. Carlos Iván Pinzón Romero}
\duedate{18 de noviembre de 2026}

% Atajos de escritura
\newcommand{\cod}[1]{\texttt{#1}}            % archivos, scripts y código
\newcommand{\ruta}[1]{\textit{#1}}           % rutas de menú: Archivo > Guardar
\newcommand{\tecla}[1]{\textsc{#1}}          % teclas

% Columnas de tabla con texto justificado a la izquierda
\newcolumntype{L}{>{\raggedright\arraybackslash}X}
\newcolumntype{P}[1]{>{\raggedright\arraybackslash}p{#1}}

% Ubicación de figuras y tablas: cerca del texto (los manuales usan H: exactamente en su lugar)
\usepackage{float}
\floatplacement{figure}{!htbp}
\floatplacement{table}{!htbp}
% Nota de figuras y tablas en español y alineada a la izquierda (APA)
\AtBeginDocument{\renewcommand{\notelabel}{Nota}}
\newcommand{\NotaFig}[1]{\if\relax\detokenize{#1}\relax\else\par\raggedright\figurenote{#1}\fi}

% Figura APA: número y título arriba, imagen y nota opcional abajo
%   \Figura[ancho]{archivo}{etiqueta}{título}{nota}
\newcommand{\Figura}[5][0.82\linewidth]{%
  \begin{figure}
    \caption{#4}\label{#3}
    \centering\includegraphics[width=#1]{#2}
    \NotaFig{#5}
  \end{figure}}

% Dos imágenes lado a lado en una misma figura
%   \FiguraDoble{archivo a}{archivo b}{etiqueta}{título}{nota}
\newcommand{\FiguraDoble}[5]{%
  \begin{figure}
    \caption{#4}\label{#3}
    \centering
    \includegraphics[width=0.49\linewidth]{#1}\hfill\includegraphics[width=0.49\linewidth]{#2}
    \NotaFig{#5}
  \end{figure}}

% Tabla APA: número y título arriba; contenido a espacio sencillo y en letra un poco menor
\newenvironment{TablaAPA}[2]{%
  \begin{table}\caption{#2}\label{#1}\centering\small\linespread{1}\selectfont}
  {\end{table}}

% Términos de las referencias como en el manual de normas APA adjunto (Menéndez Echavarría, 2020)
\DefineBibliographyStrings{spanish}{version = {Versión}, nodate = {s.\,f\adddot}}
 justo después de \documentclass.
\usepackage[spanish,es-noshorthands,es-tabla]{babel}
\usepackage[T1]{fontenc}
\usepackage{newtxtext,newtxmath}
\usepackage{csquotes}
\usepackage[style=apa,backend=biber]{biblatex}
\addbibresource{../comun/referencias.bib}
\usepackage{graphicx}
\usepackage{booktabs}
\usepackage{tabularx}
\usepackage{array}
\usepackage{enumitem}
\usepackage{xurl}
\usepackage{tikz}
\usetikzlibrary{positioning,arrows.meta,shapes.misc}
\hypersetup{hidelinks}
\setlist{itemsep=2pt,topsep=4pt,parsep=0pt}
\setlength{\emergencystretch}{3em}

% Datos de la portada
\authorsnames{Neyder Sebastián Orozco Villamil, Deisy Lorena Jiménez Martínez}
\authorsaffiliations{Universidad Central}
\course{Modelado 3D y Video Juegos (43391312)}
\professor{Mg. Carlos Iván Pinzón Romero}
\duedate{18 de noviembre de 2026}

% Atajos de escritura
\newcommand{\cod}[1]{\texttt{#1}}            % archivos, scripts y código
\newcommand{\ruta}[1]{\textit{#1}}           % rutas de menú: Archivo > Guardar
\newcommand{\tecla}[1]{\textsc{#1}}          % teclas

% Columnas de tabla con texto justificado a la izquierda
\newcolumntype{L}{>{\raggedright\arraybackslash}X}
\newcolumntype{P}[1]{>{\raggedright\arraybackslash}p{#1}}

% Ubicación de figuras y tablas: cerca del texto (los manuales usan H: exactamente en su lugar)
\usepackage{float}
\floatplacement{figure}{!htbp}
\floatplacement{table}{!htbp}
% Nota de figuras y tablas en español y alineada a la izquierda (APA)
\AtBeginDocument{\renewcommand{\notelabel}{Nota}}
\newcommand{\NotaFig}[1]{\if\relax\detokenize{#1}\relax\else\par\raggedright\figurenote{#1}\fi}

% Figura APA: número y título arriba, imagen y nota opcional abajo
%   \Figura[ancho]{archivo}{etiqueta}{título}{nota}
\newcommand{\Figura}[5][0.82\linewidth]{%
  \begin{figure}
    \caption{#4}\label{#3}
    \centering\includegraphics[width=#1]{#2}
    \NotaFig{#5}
  \end{figure}}

% Dos imágenes lado a lado en una misma figura
%   \FiguraDoble{archivo a}{archivo b}{etiqueta}{título}{nota}
\newcommand{\FiguraDoble}[5]{%
  \begin{figure}
    \caption{#4}\label{#3}
    \centering
    \includegraphics[width=0.49\linewidth]{#1}\hfill\includegraphics[width=0.49\linewidth]{#2}
    \NotaFig{#5}
  \end{figure}}

% Tabla APA: número y título arriba; contenido a espacio sencillo y en letra un poco menor
\newenvironment{TablaAPA}[2]{%
  \begin{table}\caption{#2}\label{#1}\centering\small\linespread{1}\selectfont}
  {\end{table}}

% Términos de las referencias como en el manual de normas APA adjunto (Menéndez Echavarría, 2020)
\DefineBibliographyStrings{spanish}{version = {Versión}, nodate = {s.\,f\adddot}}
 justo después de \documentclass.
\usepackage[spanish,es-noshorthands,es-tabla]{babel}
\usepackage[T1]{fontenc}
\usepackage{newtxtext,newtxmath}
\usepackage{csquotes}
\usepackage[style=apa,backend=biber]{biblatex}
\addbibresource{../comun/referencias.bib}
\usepackage{graphicx}
\usepackage{booktabs}
\usepackage{tabularx}
\usepackage{array}
\usepackage{enumitem}
\usepackage{xurl}
\usepackage{tikz}
\usetikzlibrary{positioning,arrows.meta,shapes.misc}
\hypersetup{hidelinks}
\setlist{itemsep=2pt,topsep=4pt,parsep=0pt}
\setlength{\emergencystretch}{3em}

% Datos de la portada
\authorsnames{Neyder Sebastián Orozco Villamil, Deisy Lorena Jiménez Martínez}
\authorsaffiliations{Universidad Central}
\course{Modelado 3D y Video Juegos (43391312)}
\professor{Mg. Carlos Iván Pinzón Romero}
\duedate{18 de noviembre de 2026}

% Atajos de escritura
\newcommand{\cod}[1]{\texttt{#1}}            % archivos, scripts y código
\newcommand{\ruta}[1]{\textit{#1}}           % rutas de menú: Archivo > Guardar
\newcommand{\tecla}[1]{\textsc{#1}}          % teclas

% Columnas de tabla con texto justificado a la izquierda
\newcolumntype{L}{>{\raggedright\arraybackslash}X}
\newcolumntype{P}[1]{>{\raggedright\arraybackslash}p{#1}}

% Ubicación de figuras y tablas: cerca del texto (los manuales usan H: exactamente en su lugar)
\usepackage{float}
\floatplacement{figure}{!htbp}
\floatplacement{table}{!htbp}
% Nota de figuras y tablas en español y alineada a la izquierda (APA)
\AtBeginDocument{\renewcommand{\notelabel}{Nota}}
\newcommand{\NotaFig}[1]{\if\relax\detokenize{#1}\relax\else\par\raggedright\figurenote{#1}\fi}

% Figura APA: número y título arriba, imagen y nota opcional abajo
%   \Figura[ancho]{archivo}{etiqueta}{título}{nota}
\newcommand{\Figura}[5][0.82\linewidth]{%
  \begin{figure}
    \caption{#4}\label{#3}
    \centering\includegraphics[width=#1]{#2}
    \NotaFig{#5}
  \end{figure}}

% Dos imágenes lado a lado en una misma figura
%   \FiguraDoble{archivo a}{archivo b}{etiqueta}{título}{nota}
\newcommand{\FiguraDoble}[5]{%
  \begin{figure}
    \caption{#4}\label{#3}
    \centering
    \includegraphics[width=0.49\linewidth]{#1}\hfill\includegraphics[width=0.49\linewidth]{#2}
    \NotaFig{#5}
  \end{figure}}

% Tabla APA: número y título arriba; contenido a espacio sencillo y en letra un poco menor
\newenvironment{TablaAPA}[2]{%
  \begin{table}\caption{#2}\label{#1}\centering\small\linespread{1}\selectfont}
  {\end{table}}

% Términos de las referencias como en el manual de normas APA adjunto (Menéndez Echavarría, 2020)
\DefineBibliographyStrings{spanish}{version = {Versión}, nodate = {s.\,f\adddot}}
 justo después de \documentclass.
\usepackage[spanish,es-noshorthands,es-tabla]{babel}
\usepackage[T1]{fontenc}
\usepackage{newtxtext,newtxmath}
\usepackage{csquotes}
\usepackage[style=apa,backend=biber]{biblatex}
\addbibresource{../comun/referencias.bib}
\usepackage{graphicx}
\usepackage{booktabs}
\usepackage{tabularx}
\usepackage{array}
\usepackage{enumitem}
\usepackage{xurl}
\usepackage{tikz}
\usetikzlibrary{positioning,arrows.meta,shapes.misc}
\hypersetup{hidelinks}
\setlist{itemsep=2pt,topsep=4pt,parsep=0pt}
\setlength{\emergencystretch}{3em}

% Datos de la portada
\authorsnames{Neyder Sebastián Orozco Villamil, Deisy Lorena Jiménez Martínez}
\authorsaffiliations{Universidad Central}
\course{Modelado 3D y Video Juegos (43391312)}
\professor{Mg. Carlos Iván Pinzón Romero}
\duedate{18 de noviembre de 2026}

% Atajos de escritura
\newcommand{\cod}[1]{\texttt{#1}}            % archivos, scripts y código
\newcommand{\ruta}[1]{\textit{#1}}           % rutas de menú: Archivo > Guardar
\newcommand{\tecla}[1]{\textsc{#1}}          % teclas

% Columnas de tabla con texto justificado a la izquierda
\newcolumntype{L}{>{\raggedright\arraybackslash}X}
\newcolumntype{P}[1]{>{\raggedright\arraybackslash}p{#1}}

% Ubicación de figuras y tablas: cerca del texto (los manuales usan H: exactamente en su lugar)
\usepackage{float}
\floatplacement{figure}{!htbp}
\floatplacement{table}{!htbp}
% Nota de figuras y tablas en español y alineada a la izquierda (APA)
\AtBeginDocument{\renewcommand{\notelabel}{Nota}}
\newcommand{\NotaFig}[1]{\if\relax\detokenize{#1}\relax\else\par\raggedright\figurenote{#1}\fi}

% Figura APA: número y título arriba, imagen y nota opcional abajo
%   \Figura[ancho]{archivo}{etiqueta}{título}{nota}
\newcommand{\Figura}[5][0.82\linewidth]{%
  \begin{figure}
    \caption{#4}\label{#3}
    \centering\includegraphics[width=#1]{#2}
    \NotaFig{#5}
  \end{figure}}

% Dos imágenes lado a lado en una misma figura
%   \FiguraDoble{archivo a}{archivo b}{etiqueta}{título}{nota}
\newcommand{\FiguraDoble}[5]{%
  \begin{figure}
    \caption{#4}\label{#3}
    \centering
    \includegraphics[width=0.49\linewidth]{#1}\hfill\includegraphics[width=0.49\linewidth]{#2}
    \NotaFig{#5}
  \end{figure}}

% Tabla APA: número y título arriba; contenido a espacio sencillo y en letra un poco menor
\newenvironment{TablaAPA}[2]{%
  \begin{table}\caption{#2}\label{#1}\centering\small\linespread{1}\selectfont}
  {\end{table}}

% Términos de las referencias como en el manual de normas APA adjunto (Menéndez Echavarría, 2020)
\DefineBibliographyStrings{spanish}{version = {Versión}, nodate = {s.\,f\adddot}}
